\documentclass[11pt]{article}
\usepackage[T1]{fontenc}
\usepackage[utf8]{inputenc}
\usepackage{lmodern}
\usepackage[margin=1in]{geometry}
\usepackage{amsmath,amssymb,amsthm,mathtools}
\usepackage{booktabs,array,longtable}
\usepackage{xcolor}
\usepackage{microtype}
\usepackage{xurl}
\usepackage{hyperref}
\hypersetup{colorlinks=true,linkcolor=blue!50!black,citecolor=blue!50!black,urlcolor=blue!50!black}
\newtheorem{theorem}{Theorem}[section]
\newtheorem{proposition}[theorem]{Proposition}
\newtheorem{corollary}[theorem]{Corollary}
\newtheorem{remark}[theorem]{Remark}
\newcommand{\C}{\mathbb C}
\newcommand{\R}{\mathbb R}
\newcommand{\Prob}{\mathbb P}
\newcommand{\E}{\mathbb E}
\newcommand{\tr}{\operatorname{tr}}
\newcommand{\CN}{\mathcal{CN}}
\newcommand{\Law}{\mathcal L}
\newcommand{\pd}{\succ0}
\newcommand{\lean}[1]{{\small\nolinkurl{#1}}}
\newcommand{\review}{\par\smallskip\noindent\textbf{Human semantic review required.}\ }
\title{From Gaussian Samples to a Complex Wishart Minimum Law in Lean}
\author{}
\date{Working manuscript, 16 October 2026}

\begin{document}
\maketitle
\begin{abstract}
We present a Lean formalization of selected complex Gaussian Gram matrix laws,
with a complete actual-model chain for the smallest eigenvalue of a central
complex $2\times2$ Wishart matrix. The construction starts from independent real
standard Gaussian coordinates, fixes the unit complex variance convention, and
derives an independent Bartlett coordinate law. An explicit four-real-coordinate
chart yields the normalized Gram density on the positive definite cone.
Translation of this cone, followed by a Hermitian spectral threshold criterion,
identifies the actual minimum law as exponential with rate two. We also formalize
its positive scalar variance dependence and a weakest-mode logarithmic-rate
outage formula. Supporting developments include noncentral one-column results,
an arbitrary-row one-column small-threshold asymptotic, and fixed-subspace
Gaussian projection and independence interfaces. We distinguish these actual
probability laws from general-dimensional candidate-kernel identities and
unresolved joint spectral densities. The mathematical distribution formulas
are classical; the contribution is their explicit formal connection to an
underlying Gaussian sample model and the reusable interfaces needed for that
connection. Independent semantic review and release replication remain pending
in this working version.
\end{abstract}

\section{Introduction}
Wishart eigenvalues occur in multivariate probability and in the analysis of
multiple-antenna channels. In eigenmode transmission, an eigenvalue determines
the gain of a scalar subchannel. Jin, McKay, Gao, and Collings
\cite{jin} give formulas for ordered eigenvalue distributions of complex
noncentral Wishart matrices and apply them to communication performance.
This paper studies the formal verification of selected probability statements
underlying that analysis.

The main difficulty is connecting representations. A Gaussian sample can be
represented by real coordinates, complex columns, a Gram matrix, coordinates
on the Hermitian space, and ordered eigenvalues. A density identity at one level
does not automatically establish the law at another. Variance factors, the
reference measure, exceptional sets, and strict versus weak events must be
tracked through each transition. We organize the development around equalities
of pushforward measures so that a proposed density cannot enter as an assumed
description of the model it is meant to identify.

Our principal example is a central complex matrix $H\in\C^{2\times2}$ with
independent unit-variance entries. We prove the coordinate density of $H^*H$
and derive
\[
 \Prob\{\lambda_{\min}(H^*H)>x\}=e^{-2x},\qquad x\geq0.
\]
The proof uses a translation of the positive definite cone and does not require
an ordered spectral Jacobian. Positive scalar covariance and a logarithmic-rate
statistic are then obtained by transporting events through this actual law.

The contribution has three parts. First, the central $2\times2$ chain exposes
the distributional, coordinate, and spectral obligations separately. Second,
the project provides reusable Gaussian projection and Hermitian-coordinate
interfaces, together with a dimension-independent noncentral one-column
asymptotic. Third, it supplies a declaration-level account of what has been
checked and what remains outside the formalized scope. We make no claim that
the classical distribution formulas are new, or that this is the first
formalization of Wishart theory. A comprehensive novelty comparison remains
part of release preparation.

The development builds on Lean and mathlib \cite{mathlib}. In particular,
the change-of-variables infrastructure developed in mathlib \cite{gouezel}
and its Gaussian and independence libraries carry much of the foundational
analysis. The role of probability kernels in this ecosystem is described by
Degenne \cite{degenne}. Our contribution specializes and connects those
interfaces to concrete Gaussian matrix statistics; it does not replace their
underlying general theorems.

\section{Model, measures, and formal scope}\label{sec:model}
Let $U_{ij},V_{ij}$ be independent $N(0,1)$ variables and set
\begin{equation}\label{eq:sample}
 H_{ij}=\frac{U_{ij}+\mathrm{i}V_{ij}}{\sqrt2},\qquad
 X=M+H,\qquad W=X^*X.
\end{equation}
Thus $\E|H_{ij}|^2=1$. The complex scalar density is
$\pi^{-1}\exp(-|z|^2)$ relative to $d\Re z\,d\Im z$.
The sample space in Lean is a finite-dimensional real Euclidean space named
\lean{ComplexSample}; \lean{complexSampleMatrix} performs the complex encoding.
The underlying law is mathlib's \lean{stdGaussian}. The deterministic mean
$M$ is added after encoding. The Gram law is a measurable pushforward, rather
than a measure defined initially by a Wishart density formula.

For a general $t\times s$ sample we use $X^*X$ when $s\leq t$.
For a one-column sample, the only Gram eigenvalue is $\|X\|_F^2$.
For the main example, $s=t=2$ and $M=0$. Eigenvalues are ordered decreasingly;
the last coordinate is the minimum. The formal real-valued statistic
\lean{twoRowActualMinimum} is computed from the actual sample Gram matrix.

\subsection{The Hermitian reference measure}
Write
\begin{equation}\label{eq:hermitian}
 W=\begin{pmatrix}u&b+\mathrm{i}c\\b-\mathrm{i}c&v\end{pmatrix}.
\end{equation}
The reference measure is $du\,db\,dc\,dv$. The positive definite coordinate
cone is
\[
 D=\{(u,b,c,v):u>0,\ b^2+c^2<uv\}.
\]
It is a four-dimensional real measure on Hermitian matrices. Restricting
eight-dimensional Lebesgue measure on all complex $2\times2$ matrices to the
Hermitian subspace would produce a different, unusable reference measure.

In dimension $s$, \lean{CentralHermitianCoordinates} constructs an explicit
measurable equivalence between $s^2$ real coordinates and the Hermitian
subspace. Its coordinate measure is a product of real Lebesgue measures,
transported by this equivalence. This construction fixes a normalization;
it does not by itself identify any Gaussian Gram density.

\subsection{Meaning of the verification claims}
An actual-law theorem starts with the sample measure in
\eqref{eq:sample} and identifies its pushforward. A coordinate theorem gives
an equivalence, a Jacobian, or a transformation of integrals. A candidate
identity simplifies a proposed analytical expression. These are recorded
separately. In particular, a theorem about a candidate integral does not count
as a probability theorem until its density has been derived from the sample.

The named project declarations discussed below have kernel-check evidence.
The repository's printed axiom checks use only
\lean{propext}, \lean{Classical.choice}, and \lean{Quot.sound} for the audited
declarations. This is a dependency statement, not a proof that the mathematical
interpretation of every definition matches an external paper. Independent
review of model conventions and analytical specifications remains pending.
\review The Gaussian encoding, the Hermitian coordinate measure, and the
identification with the communication model require independent scrutiny.

\section{Actual Bartlett coordinates and the Gram density}\label{sec:density}
Let $h_1,h_2$ be the columns of the central $2\times2$ sample.
Set $S=\|h_1\|^2$. On the event $h_1\ne0$, normalize the first column, project
the second column onto it, and write the resulting complex coefficient as
$Z=A+\mathrm{i}B$. Let $E$ be the squared norm of the orthogonal residual.
The implementation gives a measurable unit-direction fallback at zero.
It subsequently proves that the exceptional first-column event has mass zero.

\begin{theorem}[Actual Bartlett law]\label{thm:bartlett}
The sample-defined coordinates satisfy
\[
 \Law(S,Z,E)=\operatorname{Gamma}(2,1)\otimes\CN(0,1)
                       \otimes\operatorname{Exp}(1),
\]
where Gamma and exponential parameters are rates. Relative to
$ds\,da\,db\,de$, their density is
\begin{equation}\label{eq:bartlett}
 f_B(s,a,b,e)=\frac{s}{\pi}
           e^{-(s+a^2+b^2+e)}\,\mathbf1_{\{s>0,e>0\}}.
\end{equation}
\end{theorem}
\begin{proof}
For each fixed encoded first column, the projection and residual calculation
gives the same $\CN(0,1)\otimes\operatorname{Exp}(1)$ fiber law.
The direction map is measurable and defined everywhere. The established
Gaussian skew-product law transports this fixed-fiber statement to the actual
two-column sample, showing independence from the first column. Its energy
has the already established $\operatorname{Gamma}(2,1)$ law. Mapping this joint
law to $(S,Z,E)$ gives the displayed product measure.
The densities $s e^{-s}$, $\pi^{-1}e^{-(a^2+b^2)}$, and $e^{-e}$ multiply to
\eqref{eq:bartlett}. The positive-domain restriction is derived from these
actual marginal laws, including the absence of boundary mass.
\end{proof}

This step is stronger than establishing independence for a fixed projection.
The first-column direction is random, so measurability and the invariant fiber
law must be used explicitly. The corresponding principal declarations are
\lean{actualTwoRow_BartlettCoordinates_law} and
\lean{actualTwoRow_BartlettRealCoordinates_eq_normalizedDensity}.

The Gram coordinate chart is
\begin{equation}\label{eq:chart}
 T(s,a,b,e)=(s,\sqrt{s}\,a,\sqrt{s}\,b,a^2+b^2+e).
\end{equation}
Its domain is $s>0,e>0$, and its image is $D$. The inverse on $D$ is
\[
 (u,b,c,v)\longmapsto
 \left(u,\frac b{\sqrt u},\frac c{\sqrt u},
             v-\frac{b^2+c^2}{u}\right).
\]
The formal chart is proved to reproduce the actual Bartlett Gram mapping.
It is not an unrelated parameterization with an assumed connection to $H$.

\begin{proposition}[Jacobian and actual Gram density]\label{thm:gram}
On the positive Bartlett domain, $|\det DT|=s$. The coordinate law of $H^*H$
has density
\begin{equation}\label{eq:gram}
 f_W(u,b,c,v)=\pi^{-1}e^{-(u+v)}\mathbf1_D(u,b,c,v).
\end{equation}
Its integral is one.
\end{proposition}
\begin{proof}
The derivative of \eqref{eq:chart}, in the stated coordinate order, is
\[
 \begin{pmatrix}
 1&0&0&0\\
 a/(2\sqrt s)&\sqrt s&0&0\\
 b/(2\sqrt s)&0&\sqrt s&0\\
 0&2a&2b&1
 \end{pmatrix},
\]
whose determinant is $s$. Lean verifies the Fr\'echet derivative, the determinant,
the inverse, and the image of the domain. The change-of-variables theorem
therefore divides the Bartlett density by $s$. The factor $s$ in
\eqref{eq:bartlett} cancels, and the exponent becomes $-(u+v)$.
The actual-coordinate pushforward and its almost-sure positive support identify
the resulting measure with the actual Gram law. Normalization follows from
that law being a probability measure; it is not an additional density premise.
\end{proof}

The main declaration is
\lean{actualTwoRow_GramRealCoordinates_eq_traceDensity}; the existing
Wishart-measure API is exposed by
\lean{complexWishart_twoRow_realCoordinates_eq_traceDensity}.
\review Product-measure transport, complex-to-real volume normalization,
the orientation of the Gram coordinates, and the inverse Jacobian are retained
as explicit items in the independent review packet.

\section{From the positive cone to the minimum law}\label{sec:minimum}
For a Hermitian matrix $W$,
\begin{equation}\label{eq:threshold}
 \lambda_{\min}(W)>x\quad\Longleftrightarrow\quad W-xI\pd.
\end{equation}
The project proves this spectral criterion and its concrete four-coordinate
version. It uses the smallest coordinate of the decreasing ordered spectrum.

\begin{theorem}[Central complex $2\times2$ minimum]\label{thm:minimum}
For the actual sample \eqref{eq:sample} with $M=0$ and $s=t=2$, and $x\geq0$,
\begin{align}
 \Prob\{\lambda_{\min}(H^*H)>x\}&=e^{-2x},\label{eq:tail}\\
 \Prob\{\lambda_{\min}(H^*H)\leq x\}&=1-e^{-2x}.\label{eq:cdf}
\end{align}
For $x<0$ the CDF is zero. The complete minimum law is
$\operatorname{Exp}(2)$.
\end{theorem}
\begin{proof}
By \eqref{eq:threshold} and Proposition~\ref{thm:gram}, the tail is the integral
of $\pi^{-1}e^{-\tr W}$ over $W\succ xI$. For $x\geq0$ this region lies in
the positive definite support. Substitute $W=V+xI$. Translation preserves
the four-real-coordinate measure, while
$\tr(V+xI)=\tr V+2x$. Hence the integral is
\[
 e^{-2x}\int_{V\pd}\pi^{-1}e^{-\tr V}\,dV=e^{-2x},
\]
using the normalization already derived from the actual Gram law.
The complement of the strict tail event is exactly the weak CDF event;
no exchange of strict and weak inequalities is needed here.
Positive semidefiniteness of a Gram matrix gives the negative-threshold case.
Finally, uniqueness of probability measures with a common CDF identifies the
whole pushforward measure with $\operatorname{Exp}(2)$.
\end{proof}

The three principal statements are \lean{actualTwoRow_minimum_tail},
\lean{actualTwoRow_minimum_CDF}, and
\lean{actualTwoRow_minimum_law_expTwo}.
The unmodified determinant candidate of Jin et al.\ at
$(s,t,L)=(2,2,0)$ is also shown equal to this actual CDF in
\lean{theorem1CentralTwoRow_candidate_eq_actualCDF}.
This last connection is a specialization of their formula, not a proof of
their general ordered eigenvalue distribution theorem.

\subsection{Positive scalar covariance}
\begin{corollary}[Scale convention]\label{thm:scale}
Let $a>0$, $v=a^2$, and $W_a=(aH)^*(aH)$. Then
\[
 \lambda_{\min}(W_a)=a^2\lambda_{\min}(H^*H),\qquad
 \Law(\lambda_{\min}(W_a))=\operatorname{Exp}(2/a^2).
\]
In particular, the CDF for variance $v>0$ is $1-e^{-2x/v}$ for $x\geq0$.
\end{corollary}
\begin{proof}
The actual encoded sample matrix respects real scalar multiplication.
Consequently its Gram matrix scales by $a^2$. The spectral threshold criterion
identifies the corresponding minimum scaling, and transporting the events in
Theorem~\ref{thm:minimum} gives the CDF. CDF uniqueness again gives the full law.
\end{proof}
Here $a$ is an amplitude and $v$ is an entry variance. The result addresses a
positive scalar covariance; it does not address arbitrary correlated columns.

\section{An information-theoretic event interface}\label{sec:outage}
Consider the mathematical statistic
\[
 C_{\min}(H;a,\rho)=\log\bigl(1+\rho\lambda_{\min}((aH)^*(aH))\bigr),
 \qquad a>0,\quad\rho>0,
\]
with natural logarithms. Under a scalar Gaussian eigenmode model this is the
usual per-mode logarithmic-rate expression, measured in nats; the communication
background is described by Tse and Viswanath \cite{tse}.
The formal theorem concerns this explicitly defined statistic. It does not
formalize a coding theorem, optimal power allocation, or the total MIMO
mutual information $\log\det(I+\rho H^*H)$.

\begin{corollary}[Weakest-mode log-rate outage]\label{thm:outage}
For $R\geq0$,
\begin{equation}\label{eq:outage}
 \Prob\{C_{\min}(H;a,\rho)\leq R\}
 =1-\exp\left(-\frac{2(e^R-1)}{\rho a^2}\right).
\end{equation}
For $R<0$, the probability is zero.
\end{corollary}
\begin{proof}
For every $g\geq0$ and $\rho>0$, positivity of $1+\rho g$ and monotonicity of
the exponential give
\[
 \log(1+\rho g)\leq R\quad\Longleftrightarrow\quad
 g\leq\frac{e^R-1}{\rho}.
\]
For $R\geq0$ the transformed threshold is nonnegative, so the actual scaled
CDF applies. For $R<0$, $C_{\min}\geq0$ pointwise makes the event empty.
\end{proof}
The actual probability statement is
\lean{actualTwoRow_weakestLogRate_outage}. Its classical interpretation is
deliberately separated from the proved event transformation. A rate specified
in bits uses $R\log2$ in place of $R$; the present Lean statement uses nats.
\review The operational interpretation of $\rho$ and the allocation of power
among modes must be specified before identifying this statistic with a
particular physical outage experiment.

\section{Reusable probability foundations and supporting results}
\subsection{Fixed orthogonal Gaussian decompositions}
Let $G$ have the standard Gaussian law on a finite-dimensional real Hilbert
space $E$, and let $K$ be a fixed linear subspace. Write $P_K$ and $P_{K^\perp}$
for the subspace-valued orthogonal projections.

\begin{proposition}[Actual orthogonal split]\label{thm:split}
The projected variables have their subspace standard Gaussian laws and satisfy
\[
 \Law(P_KG,P_{K^\perp}G)=\gamma_K\otimes\gamma_{K^\perp}.
\]
If $\dim_\R K=2k$ with $k>0$, then
$\|P_KG\|^2/2\sim\operatorname{Gamma}(k,1)$.
\end{proposition}
\begin{proof}
The characteristic function of a linear image is obtained by composing with
the adjoint. For a coisometry the adjoint preserves norm, so the standard
Gaussian characteristic function is unchanged. This identifies the projected
laws. Joint Gaussianity follows by linear transport, and the cross covariance
vanishes by orthogonality. The existing Gaussian independence theorem then
gives independence and the product law. The norm statement transports the
established Euclidean radial Gamma law through an orthonormal coordinate
isometry; it does not require a new radial integration proof.
\end{proof}
These interfaces are
\lean{stdGaussian_orthogonalProjectionOnto_law},
\lean{stdGaussian_orthogonal_split_law}, and
\lean{stdGaussian_projected_half_norm_sq_law}.
They reuse mathlib's Gaussian characterization and independence machinery.
The word \emph{fixed} is essential: substituting a random earlier-column
subspace requires a separate measurable conditional-law argument.

\subsection{Noncentral one-column small-threshold behavior}
\begin{theorem}[Arbitrary positive row count]\label{thm:onecolumn}
For an integer $m>0$, a deterministic $M\in\C^{m\times1}$, and the sample
convention \eqref{eq:sample}, let
$F_M(x)=\Prob\{\|M+H\|_F^2\leq x\}$. Then
\begin{equation}\label{eq:hardedge}
 \lim_{x\downarrow0}\frac{F_M(x)}{x^m}
       =\frac{e^{-\|M\|_F^2}}{m!}.
\end{equation}
\end{theorem}
\begin{proof}
Let $\mu$ be the real encoded mean, so $\|\mu\|^2=2\|M\|_F^2$.
The actual event is a real Gaussian ball of radius $r=\sqrt{2x}$.
For its centered density $\varphi$, translation gives
\[
 \varphi(z-\mu)=e^{-\|\mu\|^2/2+\langle z,\mu\rangle}\varphi(z).
\]
Cauchy--Schwarz bounds the inner product on that ball by
$\pm r\|\mu\|$. Integrating yields
\[
 e^{-\|M\|_F^2-r\|\mu\|}F_0(x)
 \leq F_M(x)\leq
 e^{-\|M\|_F^2+r\|\mu\|}F_0(x).
\]
The central energy has the already proved $\operatorname{Gamma}(m,1)$ law,
with $F_0(x)/x^m\to1/m!$. Dividing the inequalities by $x^m$ and applying
the squeeze theorem proves \eqref{eq:hardedge}.
\end{proof}
The theorem is
\lean{noncentralOneColumnSmallestCDF_div_tendsto_frobenius}.
It includes zero mean and does not assume a noncentral spectral density.
Exact actual-model formula connections for one-column row counts one, two,
and three are also present, with the noncentrality parameter matched to the
Frobenius mean norm. They are supporting coverage, not a general multi-column
noncentral result.

\section{Verification, diagnostics, and reproducibility}\label{sec:verification}
\subsection{Declaration-level evidence}
The pinned toolchain is Lean \lean{v4.35.0-rc2}; the mathlib revision is
\lean{e18afff334217588afc719a8b707b00d90521064}.
The repository's \lean{lean-toolchain} and \lean{lake-manifest.json} fix these
dependencies. No toolchain upgrade is part of the reported development.
Table~\ref{tab:map} gives module entry points; the companion claims manifest
records exact declarations and source hashes.

\begin{table}[ht]
\centering\small
\begin{tabular}{p{0.30\linewidth}p{0.62\linewidth}}
\toprule Result & Principal source module(s)\\\midrule
Actual coordinate density & \lean{TwoRowActualDensity}; \lean{TwoRowActualNormalizedDensity}\\
Spectral tail and full CDF & \lean{TwoRowMinimumLaw}\\
Original central candidate & \lean{Theorem1CentralTwoRowActualCDF}\\
Positive scalar covariance & \lean{TwoRowScaledMinimum}\\
Log-rate outage & \lean{TwoRowLogRateOutage}\\
Projection and independent split & \lean{GaussianOrthogonalProjectionLaw}; \lean{GaussianOrthogonalSplitLaw}\\
Noncentral one-column asymptotic & \lean{T4NoncentralOneColumnAnyRows}\\
General coordinate cone integral & \lean{CentralHermitianCoordinates}; \lean{CentralWishartConeIntegral}\\
\bottomrule
\end{tabular}
\caption{Entry points for the reported statements. The last row is a
candidate-kernel identity, not an actual general Wishart law.}\label{tab:map}
\end{table}

Three verification levels are kept distinct: elaborating a leaf module,
elaborating the aggregate project entry point, and rebuilding all project
sources into fresh output. A cached entry-point check is not a clean rebuild.
The publication tooling records fresh-source re-elaboration separately from
a full Lake build and does not claim to rebuild mathlib from source.
All 164 project modules, counting generated submodules and the aggregate entry,
passed this fresh-output source check. One stale theorem identifier in an
unreferenced four-row support file was repaired without changing its statement;
the interrupted run was then resumed using only unchanged successful outputs.
This is one completed source-check campaign, not two independent reproductions.
The executable status is in the accompanying verification report; a release
must retain that report together with its exact source hashes.

The source audit removes nested comments and string literals before searching
for prohibited proof placeholders or new axiom declarations. Lexical absence
alone is insufficient: declaration-level axiom dependencies and reproducible
kernel checking provide the stronger evidence. Independent semantic review
is an additional obligation, particularly for the correspondence between
formal objects and the intended probabilistic model.

\subsection{Numerical controls}
Numerical calculations are used to reject specification mistakes before
formal proof work, and as controls for the computational interpretations.
They are not proof substitutes. The new log-rate preflight uses seed
20261015 and 200000 actual independent complex Gaussian $2\times2$ samples.
It checks the trace against $\operatorname{Gamma}(4,1)$, the minimum against
$\operatorname{Exp}(2)$, and the Hermitian quadratic eigenvalue formula against
a numerical eigensolver. The corresponding Kolmogorov--Smirnov statistics are
approximately $0.00238$ and $0.00124$; the maximum eigenvalue-formula discrepancy
is below $5\cdot10^{-15}$.

Fifteen combinations of amplitude, per-mode coefficient, and target rate
verify the event transformation and compare its empirical probability with
\eqref{eq:outage}. The reported tolerance is eight binomial standard errors
plus $8/N$, not a formal simultaneous confidence guarantee. The script and
its JSON output preserve the seed, model, controls, and individual results.

\subsection{A bounded normalization diagnostic}
A separate audit examines the project's transcription of a partition formula
associated with Theorem~3 of \cite{jin}. A Lean witness proves that the current
transcribed central candidate has real part greater than one at
$(s,t,L,k,x)=(2,3,0,2,2)$, so it cannot be an event probability at those arguments.
This is an exact statement about a specified expression in the repository.

Independent rank-one experiments at $(s,t,L)=(3,4,1)$ use seed 20261015 and
200000 actual samples for each of two noncentrality values. At noncentrality
one and threshold one, the empirical partition probability is $0.009510$;
the displayed-expression implementation gives $0.019397228$, while division
by the residual factor $\Gamma_2(3)=2$ gives $0.009698614$.
Bessel and hypergeometric implementations agree, and trace-law controls pass.
Rare-event thresholds have limited Monte Carlo resolution.

These observations do not establish an erratum to the published paper.
The journal version, possible conventions, and the exact intended ordered
domain still require source reconciliation. We preserve the original
candidate and distinguish displayed and derived normalizations. No general
Theorem~3 probability bridge is claimed. This diagnostic is ancillary to the
main minimum-law contribution.

\subsection{AI assistance and independent review}
The development used Codex for proof decomposition, implementation, numerical
controls, and drafting. The local Ollama model
\lean{deepseek-prover-v2:7b} was used as a bounded proof-completion tool for
short, precisely stated lemmas. Accepted generated proof bodies retain their
provenance and the label \lean{GUIDED_COLLABORATION} where applicable.
This workflow is not a benchmark: attempts were not sampled under a controlled
comparison protocol, and no general model-performance conclusion is drawn.

Local prover responses did not supply an adequate independent adversarial
review. A review packet identifies measure changes, Fubini uses, parameter
matching, and strict/weak spectral events for review by a different model
family and a mathematically qualified human. No such completed review is
claimed in this version. AI tools are not listed as authors; the final human
authors must verify and take responsibility for the mathematical and
bibliographic content, consistently with arXiv's tool-disclosure policy
\cite{arxivpolicy}.

\section{Generalization boundary and research direction}\label{sec:future}
For $0<s\leq t$, the desired actual central Gram density relative to a fixed
$s^2$-real-coordinate Hermitian volume is
\begin{equation}\label{eq:general}
 K_{s,t}\det(W)^{t-s}e^{-\tr W}\mathbf1_{W\pd},\qquad
 K_{s,t}=\left(\pi^{s(s-1)/2}\prod_{j=1}^s(t-j)!\right)^{-1}.
\end{equation}
The project has \emph{not} derived \eqref{eq:general} from the actual Gaussian
sample in arbitrary dimension. What is formalized is the translation identity
for the corresponding candidate-kernel integral. If the actual density and
normalization are subsequently established, the translation would give, for
$x\geq0$,
\[
 \Prob\{\lambda_{\min}>x\}
 =e^{-sx}K_{s,t}\int_{V\pd}
                \det(V+xI)^{t-s}e^{-\tr V}\,dV.
\]
The determinant factor must be retained when $t>s$. The kernel identity itself
can be stated for negative $x$; its interpretation as a tail probability then
requires the original positive-support indicator and cannot simply be reused
unchanged.

The next substantive step is a general Cholesky density induction, rather than
a collection of further low-dimensional charts. Four obligations are explicit.
First, measurable parameterized projections must lift the fixed-subspace laws
to the random earlier-column conditional setting. Second, for $r>0$ and
$z\in\C^{s-1}$ the block map
\[
 (r,z,B)\longmapsto
 \begin{pmatrix}r&\sqrt r\,z^*\\\sqrt r\,z&B+zz^*\end{pmatrix}
\]
must have a dimension-independent real-coordinate absolute Jacobian
$r^{s-1}$. Small-dimensional symbolic controls are available but are not a
general Lean proof. Third, the determinant and trace identities must combine
with this Jacobian and the recurrence
\[
 K_{s,t}=\frac{K_{s-1,t-1}}{\pi^{s-1}(t-1)!}.
\]
Finally, exceptional sets and full-rank support must be connected to the
actual pushforward measure. These obligations identify reusable probability
and multivariate-calculus contributions beyond the present example.

General noncentral ordered spectral densities and the full parameter range of
the original Theorems~1--4 remain outside the established results. In
particular, the present cone method does not supply a noncentral unitary
integration formula. The current work supplies a verified example and
interfaces for pursuing those questions, rather than a completed
formalization of the whole source paper.

\section{Conclusion}
The development connects real Gaussian samples to an actual complex Gram
density and then to a smallest-eigenvalue law, with scalar covariance and
logarithmic-rate event transport made explicit. Its main reusable lesson is
the separation of actual pushforward laws, coordinate transformations, and
candidate formulas. The arbitrary-row one-column asymptotic and fixed
orthogonal Gaussian split provide further foundations for a general density
induction. Release preparation must still complete independent semantic
review, the related-work assessment, author metadata, and a public immutable
artifact reference.

\appendix
\section{Artifact organization and release conditions}
The companion repository places Lean sources under
\path{theory/jin_wishart_formalization/}, numerical controls under
\path{experiments/}, and exact claims and verification reports under
\path{verification/}. This manuscript is self-contained LaTeX; its bibliography
is included below, with no external images, custom class, or local input files.
The files \path{claims.json}, \path{REPRODUCIBILITY.md}, and
\path{SUBMISSION_READINESS.md} accompany the working manuscript.

Before submission, the author line must be completed, an accessible repository
or archival artifact URL must be supplied, fresh-source evidence must be
checked against the release commit, and the independent review items must be
resolved or explicitly narrowed. A final related-work search must support
the precise contribution claims. The T3 diagnostic may be omitted if source
reconciliation is not complete. arXiv moderation is not peer review
\cite{arxivpolicy}; formatting compliance does not imply acceptance.

\begin{thebibliography}{9}
\bibitem{jin}
S. Jin, M. R. McKay, X. Gao, and I. B. Collings.
\emph{MIMO Multichannel Beamforming: SER and Outage Using New Eigenvalue
Distributions of Complex Noncentral Wishart Matrices}.
arXiv:cs/0611007v2, 2006.
\url{https://arxiv.org/abs/cs/0611007v2}.

\bibitem{mathlib}
The mathlib Community.
\emph{The Lean mathematical library}.
arXiv:1910.09336, 2019.
\url{https://arxiv.org/abs/1910.09336}.

\bibitem{gouezel}
S. Gou\"ezel.
\emph{A formalization of the change of variables formula for integrals in mathlib}.
arXiv:2207.12742, 2022.
\url{https://arxiv.org/abs/2207.12742}.

\bibitem{degenne}
R. Degenne.
\emph{Markov kernels in Mathlib's probability library}.
arXiv:2510.04070, 2025; version 2, 2026.
\url{https://arxiv.org/abs/2510.04070}.

\bibitem{tse}
D. Tse and P. Viswanath.
\emph{Fundamentals of Wireless Communication}.
Cambridge University Press, 2005.
Author-hosted material: \url{https://stanford.edu/~dntse/wireless_book.html}.

\bibitem{arxivpolicy}
arXiv. \emph{Content moderation}, including the policy on generative AI tools.
Accessed 15 October 2026.
\url{https://info.arxiv.org/help/moderation/index.html}.
\end{thebibliography}
\end{document}
